\uuid{YfSD}
\titre{Résolution d'équation différentielle du premier ordre}
\auteur{bodin ; croizat ; sacre}
\organisation{exo7}
\licence{CC BY-NC-SA 4.0}
\datecreate{2022-03-08}
\niveau{L1}
\module{Analyse}
\chapitre{Equation différentielle}
\sousChapitre{Résolution d'équation différentielle du premier ordre}
\difficulte{1}
\theme{}
\exo7qcmid{hilisit.equations-differentielles.040}

\contenu{
\begin{question}
\qclassification{225.01}
Soit $f(x) = e^x+3$.
De quelle(s) équations(s) différentielle(s) cette fonction est-elle solution ?
\begin{answers}
  \bad{$y' - y = e^x$}
  \good{$y' = y -3$}
  \bad{$3y'-y=0$}
  \bad{$y'-3y=0$}
\end{answers}
\begin{explanations}
Lorsqu'on dérive $f$, on obtient $f'(x) = e^x = (e^x+3)-3 = f(x) - 3$ : ainsi $f$ est solution de l'équation différentielle $y' = y - 3$. On vérifie en remplaçant dans les autres équations différentielles $y$ par $f$ (et $y'$ par $f'$) que les égalités ne sont pas vérifiées, donc que $f$ n'est pas une solution.
\end{explanations}
\end{question}
}
